% ch05 -- Das p-Bit
\label{ch:pbit}
% Quelldokument: 184

\epigraph{Das thermische Rauschen des p-Bits ist kein fundamentaler Zufall.
Es ist die makroskopisch unaufgelöste Abtastung deterministischer
Torsionskonfigurationen.}{Dok.\,184}

\section{Der Betriebspunkt zwischen Bit und Qubit}

Das klassische Bit braucht eine tiefe Potentialmulde $U \gg k_B T$ ---
robust über Jahre, ohne Kühlung. Das Qubit operiert am entgegengesetzten
Extrem: metastabile Torsion, Kohärenz im Millisekunden-Bereich, erzwungen
durch Kühlung auf mK-Temperaturen.

Das \textbf{p-Bit} besetzt den konzeptionell bislang unerschlossenen
Zwischenbereich: eine magnetische Nanostruktur, bei der die
Torsionsbarriere thermisch überwindbar wird. Das System wechselt spontan
zwischen den beiden Makrozuständen.

\section{Die natürliche Skala: \texorpdfstring{$L_8 \approx 22\,\text{nm}$}{L8 ca. 22nm}}

Die $\xi$-Längenhierarchie der FFGFT (Dok.\,173) definiert geometrisch
ausgezeichnete Skalen:
\[
  L_N \;=\; L_0 \cdot \left(\frac{1}{\xi_{\text{par}}}\right)^{\!N},
  \qquad \xi_{\text{par}} = \frac{4}{30000}
\]
Stufe $N = 8$ liefert $L_8 \approx 21{,}6\,\text{nm}$ --- gleichzeitig
exzitonischer Bohr-Radius, Transistor-Grenzgröße und Feinstruktur\-konstante.
Das ist keine Näherung, sondern geometrische Herleitung. \status{K Dok.\,184}

Auf dieser Skala ist die Torsionsbarriere durch die
Stoner-Wohlfarth-Energie $U = K_u \cdot V$ so niedrig, dass thermische
Übergänge regelmäßig stattfinden.

\begin{laierspur}
Stell dir eine Kugel auf einer flachen Bergkuppe vor. Auf einem hohen Berg
(klassisches Bit) braucht sie einen kräftigen Stoß, um die andere Seite
zu erreichen. Auf einem niedrigen Hügel (p-Bit) reicht ein Windzug ---
die normale thermische Energie der Umgebung. Die Frage ist nur:
Wie hoch muss der Hügel genau sein? FFGFT sagt: Er ist bei 22 Nanometer
von Natur aus am flachsten.
\end{laierspur}

\begin{ffgftbox}[These: Dok.\,184]
Das thermische Rauschen des p-Bits ist kein fundamentaler Zufall.
Die Übergänge folgen der sub-Planck-skaligen Geometrie auf der Zeitskala
$t_0 = t_P/7500$. Die Stochastizität ist epistemisch --- sie reflektiert
unaufgelöste Geometrie, nicht ontologischen Zufall. \status{S Dok.\,184}
\end{ffgftbox}

\section{Anwendungen und Ausblick}

p-Bits sind ein aktives Forschungsfeld in der Spintronik.
Die FFGFT-Vorhersage: p-Bit-Nanostrukturen sollten bei ${\sim}20\,\text{nm}$
ihre optimale Betriebstemperatur-zu-Barrieren-Relation finden. \status{K}

Das ist falsifizierbar: Wenn die optimale Materialskala systematisch von
22\,nm abweicht, ist die geometrische Herleitung zu überprüfen.
